\documentclass[a4paper,11pt]{ltjsarticle}
\usepackage[haranoaji]{luatexja-preset}
\usepackage{amsmath,amssymb}
\usepackage[top=12mm,bottom=10mm,left=14mm,right=14mm]{geometry}
\usepackage{array}
\usepackage{tikz}
\usetikzlibrary{calc}
\pagestyle{empty}
\setlength{\parindent}{0pt}
\newlength{\qcol}\setlength{\qcol}{0.47\textwidth}
\newlength{\qgap}
\newlength{\anscolw}\setlength{\anscolw}{0.30\textwidth}
\newlength{\ansh}
\newcommand{\Q}[2]{\parbox[t]{\qcol}{\raggedright (#1)\hspace{0.6em}$\displaystyle #2$}}
\newcommand{\QR}[4]{\Q{#1}{#2}\hfill\Q{#3}{#4}\par\vspace{\qgap}}
\newcommand{\anscell}[1]{\rule[-\ansdrop]{0pt}{\ansh}(#1)}
\newlength{\ansdrop}

\begin{document}
{\large\bfseries 計算問題　13}\hfill 名前(\underline{\hspace{32mm}})\par\vspace{3mm}
■（1）〜（16）の計算をしなさい。（17）〜（20）は方程式を解きなさい。\par\vspace{3mm}
\setlength{\qgap}{5.4mm}
\QR{1}{64x\div(-4)}{2}{(-9)\times13a}
\QR{3}{-11+(-4)}{4}{(-90)\div(+6)}
\QR{5}{(-3x+5)+(4x-1)}{6}{(-5)-(-9)+(-3)}
\QR{7}{4-(-2)\div\left(-\dfrac{2}{3}\right)}{8}{4\div3\times(-15)}
\QR{9}{(-15)\times(-13)}{10}{(24a-15)\div3}
\QR{11}{\dfrac{6}{7}\times\left\{-\dfrac{1}{3}+\left(-\dfrac{5}{6}\right)\right\}}{12}{\dfrac{1}{5}-\dfrac{1}{2}+(-0.4)}
\QR{13}{(-2)^3\times(-1)^2}{14}{\dfrac{1}{2}(6x-8)-\dfrac{1}{3}(3x-9)}
\QR{15}{2\times(-12)-18\div(-2)}{16}{(-19)-(-24)}
\QR{17}{0.6x-0.2=-2}{18}{\dfrac{1}{4}x+5=\dfrac{1}{2}x+3}
\QR{19}{3(x-5)=2(x-7)}{20}{-2x+9=-5}
\vspace{1mm}
\Q{21}{2^{10}}\par\vspace{3mm}
\setlength{\ansh}{9.5mm}\setlength{\ansdrop}{6mm}%
\begin{center}\begin{tabular}{|p{\anscolw}|p{\anscolw}|p{\anscolw}|}\hline
\anscell{1} & \anscell{2} & \anscell{3} \\\hline
\anscell{4} & \anscell{5} & \anscell{6} \\\hline
\anscell{7} & \anscell{8} & \anscell{9} \\\hline
\anscell{10} & \anscell{11} & \anscell{12} \\\hline
\anscell{13} & \anscell{14} & \anscell{15} \\\hline
\anscell{16} & \anscell{17} & \anscell{18} \\\hline
\anscell{19} & \anscell{20} & \anscell{21} \\\hline
\end{tabular}\end{center}

\end{document}
